% Part: sets-functions-relations
% Chapter: sets
% Section: unions-and-intersections

\documentclass[../../../include/open-logic-section]{subfiles}

\begin{document}

\olfileid{sfr}{set}{uni}
\olsection{اتحاد تے اشتراک}

\begin{explain}
\olref[sfr][set][bas]{sec} وچ اسی خاصیت راہیں سیٹاں دی تعریف متعارف کرائی
سی، یعنی $\Setabs{x}{\phi(x)}$ دی صورت والیاں تعریفاں۔ ایتھے اسی کوئی
خاصیت~$\phi$ ورتدے آں، تے ایس خاصیت وچ اوہناں سیٹاں دا ذکر ہو سکدا اے
جنہاں دی تعریف اسی پہلاں کر چکے آں۔ مثال لئی، جے $A$ تے~$B$ سیٹ نیں،
تاں سیٹ $\Setabs{x}{x \in A \lor x \in B}$ اوہناں ساریاں شےواں دا بنیا اے
جو $A$ یا~$B$ دے !!{element}s نیں؛ یعنی ایہ اوہ سیٹ اے جو $A$ تے~$B$ دے
\pnboblique{!!{element}s} نوں اکٹھے کردا اے۔ اسی ایہنوں \olref{fig:union}
وانگ شکل وچ ویکھ سکدے آں، جتھے اُبھاریا ہویا علاقہ دوہاں سیٹاں $A$ تے~$B$
دے \pnboblique{!!{element}s} نوں اکٹھے ظاہر کردا اے۔

\begin{figure}
  \olasset{assets/diagrams/union.tikz}
  \caption{دو سیٹاں دا اتحاد $A \cup B$، $A$ تے~$B$ دے
   \pnboblique{!!{element}s} نوں اکٹھے لے کے بنن والا سیٹ اے۔}
  \ollabel{fig:union}
\end{figure}

سیٹاں اُتے ایہ عمل---اوہناں نوں اکٹھا کرنا---بہت فائدےمند تے عام اے، ایس
لئی اسی ایہنوں اک رسمی ناں تے اک علامت دیندے آں۔
\end{explain}

\begin{defn}[اتحاد]
دو سیٹاں $A$ تے $B$ دا \emph{اتحاد}، جہنوں $A \cup B$ لکھدے نیں،
اوہناں ساریاں شےواں دا سیٹ اے جو $A$، $B$ یا دوہاں دے !!{element}s نیں۔
\[
A \cup B = \Setabs{x}{x \in A \lor x \in B}
\]
\end{defn}

\begin{ex}
کیوں جو \pnboblique{!!{element}s} نوں وار وار گِنن دا کوئی فرق نہیں پیندا،
دو ایہو جہے سیٹاں دا اتحاد جنہاں دا !!a{element} سانجھا ہووے، اوس
!!{element} نوں صرف اک واری رکھدا اے؛ مثال لئی،
$\{ a, b, c\} \cup \{ a, 0, 1\} = \{a, b, c, 0, 1\}$۔

کِسے سیٹ تے اوہدے کِسے ذیلی سیٹ دا اتحاد بس وڈا سیٹ ای اے:
$\{a, b, c \} \cup \{a \} = \{a, b, c\}$۔

سیٹ دا خالی سیٹ نال اتحاد اوہی سیٹ اے:
$\{a, b, c \} \cup \emptyset = \{a, b, c \}$۔
\end{ex}

\begin{prob}
ثابت کرو کہ جے $A \subseteq B$ ہووے، تاں $A \cup B = B$ اے۔
\end{prob}

\begin{explain}
اسی اتحاد دے ``دوہرے ہم صورت'' عمل اُتے وی غور کر سکدے آں۔ ایہ اوہ عمل
اے جو سارے اوہناں \pnboblique{!!{element}s} دا سیٹ بناندا اے جو
$A$ دے !!{element}s وی ہون تے $B$ دے !!{element}s وی۔ ایس عمل نوں
\emph{اشتراک} آکھدے نیں، تے ایہنوں \olref{fig:intersection} وانگ وکھایا
جا سکدا اے۔
\begin{figure}
  \olasset{assets/diagrams/intersection.tikz}
  \caption{دو سیٹاں دا اشتراک $A \cap B$ اوہناں دے سانجھے
    \pnboblique{!!{element}s} دا سیٹ اے۔}
  \ollabel{fig:intersection}
\end{figure}
\end{explain}

\begin{defn}[اشتراک]
دو سیٹاں $A$ تے $B$ دا \emph{اشتراک}، جہنوں $A \cap B$ لکھدے نیں،
اوہناں ساریاں شےواں دا سیٹ اے جو $A$ تے~$B$ دوہاں دے !!{element}s نیں۔
\[
A \cap B = \Setabs{x}{x \in A \land x \in B}
\]
دو سیٹاں نوں \emph{بے سانجھ} آکھدے نیں جے اوہناں دا اشتراک خالی ہووے۔
ایس دا مطلب اے کہ اوہناں دے کوئی !!{element}s سانجھے نہیں۔
\end{defn}

\begin{ex}
جے دو سیٹاں دے کوئی !!{element}s سانجھے نہ ہون، تاں اوہناں دا اشتراک
خالی اے: $\{ a, b, c\} \cap \{ 0, 1\} = \emptyset$۔

جے دو سیٹاں دے کجھ !!{element}s سانجھے ہون، تاں اوہناں دا اشتراک اوہناں
ساریاں دا سیٹ اے: $\{a, b, c \} \cap \{a, b, d \} = \{a, b\}$۔

کِسے سیٹ دا اوہدے کِسے ذیلی سیٹ نال اشتراک بس چھوٹا سیٹ ای اے:
$\{a, b, c\} \cap \{a, b\} = \{a, b\}$۔

کِسے وی سیٹ دا خالی سیٹ نال اشتراک خالی اے:
$\{a, b, c \} \cap \emptyset = \emptyset$۔
\end{ex}

\begin{prob}
پوری منطقی پابندی نال ثابت کرو کہ جے $A \subseteq B$ ہووے، تاں
$A \cap B = A$ اے۔
\end{prob}

\begin{explain}
اسی دو توں ودھ سیٹاں دا اتحاد یا اشتراک وی بنا سکدے آں۔ عام صورت وچ ایس
لئی اک سوہنا طریقہ ایہ اے: منّ لو تسیں اوہناں سارے سیٹاں نوں اکو سیٹ
وچ اکٹھا کر لیندے او جنہاں دا اتحاد (یا اشتراک) بنانا اے۔ فیر اسی اپنے
اصل سیٹاں دے اتحاد نوں اوہناں ساریاں شےواں دے سیٹ دے طور تے متعین کر
سکدے آں جو اکٹھے کیتے سیٹ دے گھٹ توں گھٹ اک !!{element} وچ ہون، تے
اشتراک نوں اوہناں ساریاں شےواں دے سیٹ دے طور تے جو اوہدے ہر !!{element}
وچ ہون۔
\end{explain}

\begin{defn}
جے $A$ سیٹاں دا سیٹ اے، تاں $\bigcup A$، سیٹ~$A$ دے
\pnboblique{!!{element}s} دے \pnboblique{!!{element}s} دا سیٹ اے:
\begin{align*}
\bigcup A & = \Setabs{x}{x \text{، سیٹ } A\text{ دے کسے !!a{element} دا عنصر اے}},
\text{ یعنی،}\\
& = \Setabs{x}{\text{کوئی } B \in A
  \text{ ایہو جہا اے کہ } x \in B}
\end{align*}
\end{defn}

\begin{defn}
جے $A$ سیٹاں دا سیٹ اے، تاں $\bigcap A$ اوہناں شےواں دا سیٹ اے جو
$A$ دے سارے عنصراں وچ سانجھیاں نیں:
\begin{align*}
\bigcap A & = \Setabs{x}{x \text{، سیٹ } A\text{ دے ہر !!{element} وچ اے}},
\text{ یعنی،}\\
 & = \Setabs{x}{\text{ہر } B \in A, x \in B}
\end{align*}
\end{defn}

\begin{ex}
منّ لو کہ $A = \{ \{ a, b \}, \{ a, d, e \}, \{ a, d \} \}$۔
فیر $\bigcup A = \{ a, b, d, e \}$ تے $\bigcap A = \{ a \}$ اے۔
\end{ex}
\begin{prob}
	وکھاؤ کہ جے $A$ اک سیٹ اے تے $A \in B$، تاں $A \subseteq \bigcup B$ اے۔
\end{prob}

اسی سیٹاں دے سلسلے $A_1$، $A_2$، \dots لئی وی ایہو کم کر سکدے آں:
\begin{align*}
\bigcup_i A_i & = \Setabs{x}{x \text{، سیٹاں } A_i\text{ وچوں کسے اک وچ شامل اے}}\\
\bigcap_i A_i & = \Setabs{x}{x \text{، سیٹاں } A_i\text{ وچوں ہر اک وچ شامل اے}}.
\end{align*}

جدوں ساڈے کول سیٹاں دا \emph{اشاریہ} ہووے، یعنی کوئی سیٹ $I$ جس لئی
اسی $A_i$ اُتے غور کردے ہاں، ہر $i \in I$ لئی، تاں اسی ایہ مختصر
لکھتاں وی ورت سکدے آں:
\begin{align*}
	\bigcup_{i \in I} A_i & = \bigcup \Setabs{A_i }{i \in I}\\
	\bigcap_{i \in I} A_i & = \bigcap\Setabs{A_i}{i \in I}
\end{align*}

آخر وچ، اسی سارے اوہناں \pnboblique{!!{element}s} دے سیٹ بارے سوچ سکدے
آں جو $A$ وچ نیں پر $B$ وچ نہیں۔ ایہنوں \olref{difference} وانگ وکھایا
جا سکدا اے۔

\begin{figure}
  \olasset{assets/diagrams/difference.tikz}
  \caption{دو سیٹاں دا فرق $A \setminus B$ اوہناں \pnboblique{!!{element}s}
    دا سیٹ اے جو $A$ وچ نیں پر $B$ دے !!{element}s نہیں۔}
  \ollabel{difference}
\end{figure}

\begin{defn}[فرق]
\emph{سیٹاں دا فرق}~$A \setminus B$ سارے اوہناں \pnboblique{!!{element}s}
دا سیٹ اے جو $A$ وچ نیں پر $B$ دے !!{element}s نہیں، یعنی
\[
A\setminus B = \Setabs{x}{x\in A \text{ تے } x \notin B}.
\]
\end{defn}

\begin{prob}
	ثابت کرو کہ جے $A \subsetneq B$ ہووے، تاں $B \setminus A \neq \emptyset$ اے۔
\end{prob}

\end{document}
